\section{Rectangular consequences of the uniform entropy bound}
\label{sec:rectangular_specialization}

Let $V=\{0,\ldots,w-1\}\times\{0,\ldots,h-1\}$, with the open
nearest-neighbor graph and site spaces $\mathbb C^q$.
Either side length may be zero. The following reduction is the
rectangular and vertex-boundary implication in
\cite[Corollary~1.2]{OpenAI2026AreaLaw}.
% The unconditional corollary remains open: the final theorem below assumes
% the uniform finite-domain area law. No new entry is marked complete before
% hosted linter-bearing validation of the precise candidate revision.

\begin{corollary}[Rectangles and vertex boundaries]
    \label{cor:rectangles}
    For fixed $q\ge1$ and $J,\Delta>0$, there is a uniform finite constant
    $C$ such that every subset of a finite rectangular square lattice with
    bounded on-site and nearest-neighbor interactions and a unique ground
    state separated by global gap $\Delta$ satisfies
    $S_\Omega(A)\le C|\partial_V A|$.
    For arbitrary range $R$, every finite induced domain covered by the
    area-law theorem also satisfies the two vertex bounds
    $S_\Omega(A)\le4C|\partial_{\rm in}A|$ and
    $S_\Omega(A)\le2C|Z_A|$, with that theorem's constant $C(q,R,J,\Delta)$.
    This is \cite[Corollary~1.2]{OpenAI2026AreaLaw}.
\end{corollary}
% No Lean ownership/completion marker: the unconditional source corollary
% awaits the source-faithful proof of the uniform finite-domain area law.

\begin{definition}[Native rectangular interaction family]
    \label{def:rectangular_hamiltonian}
    \lean{TNLean.PEPS.AreaLaw.RectangleConfiguration,
        TNLean.PEPS.AreaLaw.RectangleStateSpace,
        TNLean.PEPS.AreaLaw.RectangularHamiltonian,
        TNLean.PEPS.AreaLaw.RectangularHamiltonian.siteTerm,
        TNLean.PEPS.AreaLaw.RectangularHamiltonian.edgeTerm,
        TNLean.PEPS.AreaLaw.RectangularHamiltonian.site_supported,
        TNLean.PEPS.AreaLaw.RectangularHamiltonian.edge_supported,
        TNLean.PEPS.AreaLaw.RectangularHamiltonian.site_norm_le,
        TNLean.PEPS.AreaLaw.RectangularHamiltonian.edge_norm_le,
        TNLean.PEPS.AreaLaw.RectangularHamiltonian.hermitian,
        TNLean.PEPS.AreaLaw.RectangularHamiltonian.operator,
        TNLean.PEPS.AreaLaw.RectangularHamiltonian.IsGappedGroundState,
        TNLean.PEPS.AreaLaw.RectangleInteraction,
        TNLean.PEPS.AreaLaw.rectangleInteractionSupport}
    For each site $v$ and undirected nearest-neighbor edge $e$, choose matrices
    $h_v,h_e$, supported on $\{v\}$ and $e$, respectively, with
    $\|h_v\|,\|h_e\|\le J$. Only their sum
    $H=\sum_v h_v+\sum_e h_e$ is required to be Hermitian.
    A unit vector $\Omega$ of energy $E_0$ satisfies the global gap condition
    $H-E_0 I\ge\Delta(I-|\Omega\rangle\langle\Omega|)$.
    The labels of the summands are the disjoint union of the sites and edges;
    their designated supports are their singleton and endpoint sets.
\end{definition}

\begin{theorem}[Lossless Hermitian replacement]
    \label{thm:rectangular_hermitian_terms}
    \lean{TNLean.PEPS.AreaLaw.RectangularHamiltonian.hermitianTerm,
        TNLean.PEPS.AreaLaw.RectangularHamiltonian.hermitianTerm_mem_supportedOperators,
        TNLean.PEPS.AreaLaw.RectangularHamiltonian.hermitianTerm_isHermitian,
        TNLean.PEPS.AreaLaw.RectangularHamiltonian.hermitianTerm_norm_le,
        TNLean.PEPS.AreaLaw.RectangularHamiltonian.sum_hermitianTerm}
    Replacing every summand $h_i$ by $k_i=(h_i+h_i^*)/2$ preserves its
    designated support, gives $k_i=k_i^*$ and $\|k_i\|\le J$, and preserves
    the total: $\sum_i k_i=H$.
\end{theorem}
\begin{proof}
    \uses{def:rectangular_hamiltonian}
    The supported matrix algebra is closed under adjoints and linear
    combinations. The Hermitian-part map is contractive and real-linear,
    so $\sum_i k_i=(H+H^*)/2=H$.
\end{proof}

\begin{theorem}[Identification with range-one domain interactions]
    \label{thm:rectangular_domain_hamiltonian}
    \lean{TNLean.PEPS.AreaLaw.rectangleAreaLawSupport,
        TNLean.PEPS.AreaLaw.rectangleAreaLawSupport_injective,
        TNLean.PEPS.AreaLaw.rectangleAreaLawState,
        TNLean.PEPS.AreaLaw.norm_rectangleAreaLawState,
        TNLean.PEPS.AreaLaw.embedOp_rectangular_eq_reindex,
        TNLean.PEPS.AreaLaw.RectangularHamiltonian.toAreaLawHamiltonian,
        TNLean.PEPS.AreaLaw.RectangularHamiltonian.toAreaLawHamiltonian_operator,
        TNLean.PEPS.AreaLaw.RectangularHamiltonian.isGappedGroundState_toAreaLawHamiltonian}
    Let $U$ identify the native configuration basis with the basis of the
    corresponding integer lattice domain. For $J\ge0$, placing $Uk_iU^*$
    at its singleton or pair support, and zero at unused admissible supports,
    gives a range-one interaction family with bound $J$ and total $UHU^*$.
    The vector $U\Omega$ has the same norm and ground energy, and satisfies
    the same global projector gap with parameter $\Delta$.
\end{theorem}
\begin{proof}
    \uses{thm:rectangular_hermitian_terms}
    Singletons and adjacent pairs have induced graph diameter at most one.
    Distinct native sites and undirected edges have distinct supports, so no
    two summands accumulate at one support. Unitary conjugation preserves
    support, adjoints, operator norms, the eigenvector equation and positive
    semidefiniteness. Summing the placed matrices therefore gives $UHU^*$.
\end{proof}

\begin{definition}[Native vertex boundaries]
    \label{def:rectangular_vertex_boundaries}
    \lean{TNLean.PEPS.AreaLaw.rectangleInnerBoundary,
        TNLean.PEPS.AreaLaw.rectangleEndpointBoundary}
    For $A\subseteq V$, the inner boundary is the set of $x\in A$ with
    a neighbor outside $A$. The set $Z_A$ consists of both endpoints of
    every edge crossing $A$.
\end{definition}

\begin{theorem}[Exact edge and vertex boundary counts]
    \label{thm:rectangular_boundary_transport}
    \lean{TNLean.PEPS.AreaLaw.card_edgeBoundary_rectangular,
        TNLean.PEPS.AreaLaw.mem_endpointBoundary_iff_adj,
        TNLean.PEPS.AreaLaw.map_rectangleInnerBoundary,
        TNLean.PEPS.AreaLaw.map_rectangleEndpointBoundary,
        TNLean.PEPS.AreaLaw.card_innerBoundary_rectangular,
        TNLean.PEPS.AreaLaw.card_endpointBoundary_rectangular,
        TNLean.PEPS.AreaLaw.rectangleEdgeBoundary_card_le_four_mul_innerBoundary_card,
        TNLean.PEPS.AreaLaw.rectangleEdgeBoundary_card_le_two_mul_endpointBoundary_card}
    The native crossing edges, inner boundary and endpoint set map
    bijectively to the corresponding integer-domain sets. Consequently,
    $|\partial_V A|\le4|\partial_{\rm in}A|$ and
    $|\partial_V A|\le2|Z_A|$, including empty boundaries.
\end{theorem}
\begin{proof}
    \uses{def:rectangular_vertex_boundaries}
    The coordinate bijection preserves adjacency and cut membership.
    Sorting the two native endpoints counts each undirected edge exactly
    once. The degree bounds on the induced domain then apply with unchanged
    cardinalities: each inner site has at most four crossing neighbors,
    while counting both endpoints counts every crossing edge twice.
\end{proof}

\begin{definition}[Uniform rectangular entropy bounds]
    \label{def:rectangular_entropy_area_law}
    \lean{TNLean.PEPS.AreaLaw.RectangularEntropyAreaLaw}
    For every integer $q\ge1$ and $J,\Delta>0$, there is $C\ge0$, chosen
    before $w,h$, the interaction family, the ground vector and the cut,
    such that all the hypotheses of Definition~\ref{def:rectangular_hamiltonian}
    imply simultaneously
    \begin{align}
        S_\Omega(A) & \le C|\partial_V A|,
                    & S_\Omega(A)          & \le4C|\partial_{\rm in}A|,
                    & S_\Omega(A)          & \le2C|Z_A|.
        \label{eq:rectangular_entropy_bounds}
    \end{align}
\end{definition}

\begin{theorem}[Rectangular bounds from the uniform finite-domain area law]
    \label{thm:rectangular_entropy_from_uniform}
    \lean{TNLean.PEPS.AreaLaw.regionalEntropy_rectangleAreaLawState,
        TNLean.PEPS.AreaLaw.UniformAreaLaw.rectangularEntropyAreaLaw}
    Assume the uniform finite-domain area law of
    \cite[Theorem~1.1]{OpenAI2026AreaLaw}. Then
    Definition~\ref{def:rectangular_entropy_area_law} holds, with the
    range-one constant $C(q,1,J,\Delta)$.
\end{theorem}
\begin{proof}
    \uses{thm:rectangular_domain_hamiltonian,thm:rectangular_boundary_transport}
    Choose the uniform constant at range one before the rectangle and
    physical instance. Apply the finite-domain theorem to $UHU^*$ and
    $U\Omega$. Relabeling tensor factors preserves the regional entropy,
    and the boundary bijections preserve all three cardinalities.
    Multiplication of the two counting inequalities by $C\ge0$ gives
    the two vertex forms in~(\ref{eq:rectangular_entropy_bounds}).
\end{proof}

\begin{theorem}[Equal-side consequence]
    \label{thm:rectangular_equal_side_entropy}
    \lean{TNLean.PEPS.Approximation.SquareHamiltonian.toRectangularHamiltonian,
        TNLean.PEPS.AreaLaw.RectangularEntropyAreaLaw.squareEntropyAreaLaw}
    The rectangular entropy bound implies the square-grid entropy statement
    used in \cite[Proposition~2.1]{OpenAI2026PolynomialPEPS}, with its original
    restriction $q\ge2$.
\end{theorem}
\begin{proof}
    \uses{def:rectangular_hamiltonian,def:rectangular_entropy_area_law}
    Put $w=h=L$, keep the original site and edge matrices, and take the
    edge-boundary component of~(\ref{eq:rectangular_entropy_bounds}).
\end{proof}
